%===================================================================================================
%                                            CAPA
%===================================================================================================
\thispagestyle{fancy}
\fancyhf{}
\lhead{
  \begin{picture}(225,150)%(horizontal,vertical)
    \put(0,0){\includegraphics[scale=0.13]{pre/logo-unicamp.jpg}}
  \end{picture}
}
\cfoot{\thepage}


\begin{center}
\vspace*{6.cm}

Gilberto Luis Valente Costa

\vspace{2.6cm}
{\fontsize{23}{23} \textbf{$hp^2$FEM - Uma Arquitetura de \textit{Software}}}	\\ \vspace{1.2ex}
{\fontsize{23}{23} \textbf{$p$ Não-Uniforme para o Método de}}			\\ \vspace{1.2ex}
{\fontsize{23}{23} \textbf{Elementos Finitos de Alta Ordem}}


\vspace*{4.6cm}
\begin{flushleft}
119/2012
% (\textcolor[rgb]{0.00,0.00,1.00}{Número de publicação - solicitar na CPG/FEM}).
\end{flushleft}

\vspace{3cm}
CAMPINAS\\ 2012
\end{center}
%===================================================================================================


%===================================================================================================
%                                     ASSINATURA DO ORIENTADOR
%===================================================================================================
%%% Força a contagem da página.
\clearpage
\setcounter{page}{3}
%%%

\thispagestyle{fancy}
\fancyhf{}
\lhead{
  \begin{picture}(225,150)%(horizontal,vertical)
    \put(0,0){\includegraphics[scale=0.13]{pre/logo-unicamp.jpg}}
  \end{picture}
}
\cfoot{\thepage}


\begin{center}
\vspace*{2.2cm}
{\large\textbf{UNIVERSIDADE ESTADUAL DE CAMPINAS\\\vspace{1.2ex}
FACULDADE DE ENGENHARIA MECÂNICA\vspace{1.2ex}}}

\vspace{1.5cm}
{\large Gilberto Luis Valente da Costa}

\vspace{1.3cm}
{\fontsize{23}{23} \textbf{$hp^2$FEM - Uma Arquitetura de \textit{Software}}}	\\ \vspace{1.2ex}
{\fontsize{23}{23} \textbf{$p$ Não-Uniforme para o Método de}}			\\ \vspace{1.2ex}
{\fontsize{23}{23} \textbf{Elementos Finitos de Alta Ordem}}
\end{center}

\vspace{0.5cm}

\begin{flushright}
\parbox{3.5in}{ Dissertação de Mestrado (ou Tese de Doutorado) apresentada à Faculdade de 
Engenharia Mecânica da Universidade Estadual de Campinas como parte dos requisitos exigidos para 
obtenção do título de Mestre(a) (ou Doutor(a)) em Engenharia Mecânica, na Área de }
\end{flushright}

\vspace{0.5cm}
\noindent
Orientador: Prof. Dr. Marco Lúcio Bittencourt\\
Coorientador:

\vspace{0.4cm}
\begin{flushleft}
\parbox{3.5in}{ESTE EXEMPLAR CORRESPONDE À VERSÃO FINAL DA DISSERTAÇÃO DEFENDIDA PELO ALUNO
Gilberto Luis Valente da Costa, E
ORIENTADO PELO PROF. DR.
Marco Lúcio Bittencourt.
\begin{center}
...............................................................\\
ASSINATURA DO ORIENTADOR
\end{center}
}
\end{flushleft}

\vspace{0.4cm}
\begin{center}
CAMPINAS\\ 2012
\end{center}
%===================================================================================================


%===================================================================================================
%                                     FICHA CATALOGRÁFICA
%===================================================================================================
%%% Força a contagem da página.
\clearpage
\setcounter{page}{4}
%%%

\begin{center}
\begin{tabular}{c}
  FICHA CATALOGRÁFICA ELABORADA PELA \\
  BIBLIOTECA DA ÁREA DE ENGENHARIA E ARQUITETURA - BAE - UNICAMP \\
\end{tabular}

\vspace{0.8cm}
\begin{tabular}{|cl|} \hline
  \hspace{2cm} & \\
  & Valente da Costa, Gilberto Luis, 1983-				\\  
  V234a & \hspace{0.15in} hp2FEM - uma arquitetura de software		\\
  & p não-uniforme para o método de elementos finitos de alta		\\
  & ordem $\backslash$ Gilberto Luis Valente				\\
  & da Costa. --Campinas, SP: [s.n.], 2012.				\\
  &  \\  
  & \hspace{0.15in} Orientador: Marco Lúcio Bittencourt. 		\\
  & \hspace{0.15in} Dissertação de Mestrado - Universidade Estadual de\\
  & Campinas, Faculdade de Engenharia Mecânica. 			\\
  & \\
  & \hspace{0.15in} 1. Arquitetura de software. 2. Método de elementos\\
  & finitos. 3. Framework (Programa de computador). I.			\\
  & Bittencourt, Marco Lúcio, 1972-. II. Universidade			\\
  & Estadual de Campinas. Faculdade de Engenharia			\\
  & Mecânica. III. hp2FEM - uma arquitetura de software		\\
  & p não-uniforme para o método de elementos finitos de alta.		\\
  & \\ \hline
\end{tabular}
\end{center}

\vspace{0.2in}
\begin{tabular}{ll}
Título em Inglês: &
hp2FEM - a p non-uniform software architecture to the high order	\\
& finite element method.						\\
Palavras-chave em Inglês: &
Software Architecture, Finite Element Method, Framework.		\\
Área de concentração: & Mecânica dos Sólidos e Projeto Mecânico	\\
Titulação: &
Mestre em Engenharia Mecânica						\\
Banca Examinadora: & Carlos Alberto Cimini Júnior, Edson Borin		\\
Data da defesa: & 27-08-2012						\\
Programa de Pós Graduação:& Engenharia Mecânica
\end{tabular}
\fancyhead{}	% Limpa os campos do cabeçalho atual
%===================================================================================================


%===================================================================================================
%                                     FOLHA DE APROVAÇÃO
%===================================================================================================
%%% Força a contagem da página.
\clearpage
\setcounter{page}{5}
%%%

\begin{center}

\vspace*{1cm}
{\large\textbf{UNIVERSIDADE ESTADUAL DE CAMPINAS\\\vspace{1.2ex}
FACULDADE DE ENGENHARIA MECÂNICA\\\vspace{1.2ex}
COMISSÃO DE PÓS-GRADUAÇÃO EM ENGENHARIA MECÂNICA\\\vspace{1.2ex}
DEPARTAMENTO DE PROJETO MECÂNICO}}

\vspace{0.8cm}
\textbf{DISSERTAÇÃO DE MESTRADO ACADÊMICO}

\vspace{1.4cm}
{\fontsize{23}{23} \textbf{$hp^2$FEM - Uma Arquitetura de \textit{Software}}}	\\ \vspace{1.4ex}
{\fontsize{23}{23} \textbf{$p$ Não-Uniforme para o Método de}}			\\ \vspace{1.4ex}
{\fontsize{23}{23} \textbf{Elementos Finitos de Alta Ordem}}			\\

\begin{flushleft}
\vspace{0.2cm}
Autor: Gilberto Luis Valente da Costa\\Orientador:~Marco~Lúcio~Bittencourt

\vspace{0.4cm}
A Banca Examinadora composta pelos membros abaixo aprovou esta Dissertação:  \\

\vspace{0.6cm}
\hrulefill\\
\textbf{Prof. Dr. Marco Lúcio Bittencourt, Presidente\\
DPM/FEM/UNICAMP \\}

\vspace{0.6cm}
\hrulefill\\
\textbf{Prof. Dr. Carlos Alberto Cimini Júnior\\
DPM/FEM/UNICAMP \\}

\vspace{0.6cm}
\hrulefill\\
\textbf{Prof. Dr. Edson Borin\\
IC/UNICAMP \\}
\end{flushleft}

\begin{flushright}
\vspace{1.2cm}
Campinas, 27 de Agosto de 2012.
\end{flushright}

\end{center}
%===================================================================================================


%===================================================================================================
%                                         DEDICATÓRIA
%===================================================================================================
%%% Força a contagem da página.
\clearpage
\setcounter{page}{7}
%%%

\begin{center}
\chapter*{Dedicatória}
\end{center}
\vspace*{1cm}

Aos meus queridos avós, Pedro Catarino e Luiza Marçal, os quais sinto muito orgulho e agradeço por
todo carinho e o apoio de sempre.
%===================================================================================================


%===================================================================================================
%                                       AGRADECIMENTOS
%===================================================================================================
%%% Força a contagem da página.
\clearpage
\setcounter{page}{9}
%%%

\begin{center}
\chapter*{Agradecimentos}
\end{center}

\vspace*{1cm}
\begin{trivlist}  \itemsep 2ex  \normalsize

\item À Deus, minha fonte de força, fé e determinação. E como em todos os dias, agradeço pelo dom da
vida.

\item Ao meu orientador, Prof. Dr. ...

\item Aos membros das bancas de qualificação e defesa ...

\item A toda minha família ...

\item Aos amigos e companheiros de laboratório ...

\item A CAPES - Coordenação de Aperfeiçoamento de Pessoal de Nível Superior, pelo indispensável
apoio financeiro.

\end{trivlist}
%===================================================================================================


%===================================================================================================
%                                         EPÍGRAFE
%===================================================================================================
%%% Força a contagem da página.
\clearpage
\setcounter{page}{11}
%%%

\vspace*{6in}
\epigraph{\emph{O gênio consiste em um por cento de inspiração e noventa e nove por cento de
transpiração.}}{Thomas Alva Edison}
%===================================================================================================


%===================================================================================================
%                                          RESUMO
%===================================================================================================
%%% Força a contagem da página.
\clearpage
\setcounter{page}{13}
%%%

\begin{center}
\chapter*{Resumo}
\end{center}
\vspace{24pt}
\onehalfspacing
\noindent
VALENTE DA COSTA, Gilberto Luis. hp2FEM - uma arquitetura de software p não-uniforme para o método
de elementos finitos de alta ordem. 2012. 122p. Dissertação (Mestrado). Faculdade de Engenharia
Mecânica, Universidade Estadual de Campinas, Campinas.\\

Este trabalho tem como objetivo principal a implementação de uma arquitetura de \textit{software}
para o Método de Elementos Finitos de Alta Ordem (MEF-AO), baseando-se no paradigma de programação
orientada a objeto (POO) e no uso de técnicas de otimização de código fonte. O \textit{software} foi
escrito em linguagem C++ e desenvolvido sobre um \textit{framework} com ferramentas que auxiliaram
no desenvolvimento. A modelagem do sistema foi realizada de forma a facilitar e promover o reuso e
manutenção do código. Buscou-se, também, a flexibilidade e generalização do MEF-AO ao permitir a
variação nos procedimentos da construção das equações e o uso de malhas $p$ não-uniforme. Neste
caso, cada elemento pode ser interpolado com uma ordem polinomial diferente, além de permitir o uso
de um algoritmo local de solução. Tal característica pode diminuir o número de operações e de
armazenamento, pois o número de funções de forma é aumentado apenas onde é necessário o uso de mais
pontos para interpolação da malha de solução. No final, o \textit{software} é avaliado aplicando o
problema de projeção para malha de quadrados e hexaedros.\\

\vspace{1cm}
\noindent
\emph{ Palavras-chave}: Arquitetura de \textit{Software}, Método de Elementos Finitos,
\textit{Framework} (Programa de computador).\\
%===================================================================================================


%===================================================================================================
%                                        ABSTRACT
%===================================================================================================
%%% Força a contagem da página.
\clearpage
\setcounter{page}{15}
%%%

\begin{center}
\chapter*{Abstract}
\end{center}
\vspace{24pt}
\noindent
VALENTE DA COSTA, Gilberto Luis. hp2FEM - a p non-uniform software architecture to the high order
finite element method. 2012. 122p. Dissertação (Mestrado). Faculdade de Engenharia Mecânica,
Universidade Estadual de Campinas, Campinas.\\

The main objective of this work is the implementation of a software architecture for the High-Order
Finite Element Method (HO-FEM), based on the Object Oriented Paradigm (OOP) and on source-code
optimization techniques. The software was written in C++ programming language and developed over a
framework which provided tools that assisted the implementation. The system was modeled so to
promote code reuse and maintainability. Furthermore, the system modeling also provided flexibility
and generalization for the HO-FEM by allowing modifications on the procedures for equation
assembling and the use of $p$-non-uniform meshes. In this case, each element can be interpolated
with different polynomial order, and allows the application of an algorithm for local solution. Such
features can reduce the number of operations for memory allocation, since the number of shape
functions is increased only where a higher density of points is needed by the solution mesh.
Finally, the software is assessed by applying the projection problem for meshes of squares and
hexahedros.\\

\vspace{1cm}
\noindent
\emph{Keywords}: Software Architecture, Finite Element Method, Framework.
%===================================================================================================


%===================================================================================================
   

%===================================================================================================
%                                    LISTA DE FIGURAS
%===================================================================================================
%%% Força a contagem da página.
\clearpage
\setcounter{page}{17}
%%%

\phantomsection
\addcontentsline{toc}{chapter}{Lista de Ilustrações}
\DeclareRobustCommand\listfigurename
{\begin{center} {\chapter*{Lista de Ilustrações}} \end{center}}

\listoffigures
%===================================================================================================


%===================================================================================================
%                                    LISTA DE TABELAS
%===================================================================================================
%%% Força a contagem da página.
\clearpage
\setcounter{page}{19}
%%%

\phantomsection
\addcontentsline{toc}{chapter}{Lista de Tabelas}
\DeclareRobustCommand\listtablename
{\begin{center} {\chapter*{Lista de Tabelas}} \end{center}}

\listoftables
%===================================================================================================


%===================================================================================================
%                                    LISTA DE SÍMBOLOS
%===================================================================================================
%%% Força a contagem da página.
\clearpage
\setcounter{page}{21}
%%%

\phantomsection
\addcontentsline{toc}{chapter}{Lista de Abreviaturas e Siglas}
\renewcommand{\listfigurename}
{\begin{center} {\chapter*{Lista de Abreviaturas e Siglas}} \end{center}}
\markboth{Lista de Abreviaturas e Siglas}{Lista de Abreviaturas e Siglas}
\printnomenclature
\chapter*{List of Symbols and Acronyms}
\addcontentsline{toc}{chapter}{List of Symbols and Acronyms}
%\vspace{-2.0cm}
\section*{Symbols}

\setlength{\LTleft}{10pt}

\subsection*{Chapter 1}

\begin{longtable}{ll}
	$\mathbb{N}^+$			& positive integers; \\
	$\mathbb{R}$			& real numbers; \\
	$\mathbb{R}^n$			& $\mathbb{R}^{nx1}$ $n$-dimensional Euclidean space; \\
	
	$\forall$				& for all; \\
	$\in$					& belongs to; \\
	$>$						& greater than; \\
	$<$						& less than; \\
	$\geq$					& greater than or equal to; \\
	$\leq$					& less than or equal to; \\
	$\triangleq$			& equals by definition; \\	
	$\approx$				& approximately equal to; \\
	$\sim$					& is distributed as; \\
	
	$\mathcal{N}$			& Gaussian distribution; \\
	$\mathcal{E}$			& Exponential distribution; \\
	$x_i$					& $i$th entry of x; \\	
	
	$f$						& process model; \\
	$g$						& observation model; \\
	$t$						& continuous-time index; \\
	$k$						& discrete-time index of measurements; \\
	$i$						& discrete-time index of input; \\
	$t_k$					& continuous-time sampled instants; \\
	$T$						& input sampling interval; \\
	$x(t)$					& state vector; \\
	$x(t_k)$				& state vector at continuous-time sampled instants; \\
	$u(t)$					& input vector; \\
	$u(iT)$					& input vector at continuous-time sampled at time $t = iT$; \\						
	$w(t)$					& process noise vector; \\
	$y_{t_k}$            	& output vector at continuous-time sampled instants;\\
	$\delta_{k}$			& measurements time-delay; \\
	$v(t_k)$				& measurement noise vector; \\
	$h_k$					& time intervals between two continuous-time sampled instants; \\
	$\lambda$ 				& parameter of the exponential distribution; \\
	$\lambda_h$				& parameter of the exponentially distributed random variable $h_k$; \\
	$\lambda_{\delta_{k}}$	& parameter of the exponentially distributed random variable $\delta_{k}$; \\
	$N$						& amount of identical sensors in the sampling model; \\
	$L$						& sampling period of the identical sensors in the sampling model; \\
	$P$						& covariance matrix; \\
	$\alpha$				& relation between output (expected) and input sampling time intervals . \\
\end{longtable}

\subsection*{Chapter 2}

\begin{longtable}{ll}
	$\mathcal{N}$			& Gaussian distribution; \\
	$\mathbb{P}$			& power set; \\
	
	$\mu$ 					& mean of a random variable; \\
	$\sigma$ 				& standard deviation of a random variable;\\
	
	$n$						& number of homogeneous sensors; \\
	$E$						& evidence; \\ 
	$H$						& hypothesis; \\
	$R_r$					& $r^{th}$ fuzzy rule; \\
	$x_k$					& $k^{th}$ input fuzzy variable; \\
	$A_i$					& $i^{th}$ antecedent linguistic value; \\
	$y$						& output variable; \\
	$C_j$					& $j^{th}$ consequent class; \\
	$A$						& fuzzy set; \\
	$U$						& universe of discourse; \\
	$B$						& subset of features; \\
	$B_*(X)$				& \textit{B-lower} approximation of $X$; \\
	$B^*(X)$				& \textit{B-upper} approximation of $X$; \\
	$BN_B(X)$				& \textit{B-boundary} region of $X$; \\
	
	$\rho(x)$				& probability distribution of $x$; \\
	$\rho(A|B)$				& conditional probability distribution of $A$ given $B$; \\
	$\pi_x(u)$				& possibility distribution of $x$ associated with $u$; \\
	$\mu_A(x)$				& membership function of $x$ of a fuzzy se $A$; \\
	
	
	$\subseteq$				& is a subset of; \\
	$\cap$					& set intersection; \\
	$\emptyset$				& empty set; \\
	
\end{longtable}


\subsection*{Chapter 3}

\begin{longtable}{ll}
	$\mathbb{N}$			& natural numbers; \\
	$\mathbb{R}$			& real numbers; \\
	$\mathbb{R}^n$			& $\mathbb{R}^{nx1}$ $n$-dimensional Euclidean space; \\
	
	$Ber(p)$				& Bernoulli distribution; \\
	$\mathcal{N}$			& Gaussian distribution; \\
	$\Gamma$				& gamma distribution; \\
	
	$\forall$				& for all; \\
	$>$						& greater than; \\
	$\geq$					& greater than or equal to; \\
	$\sim$					& is distributed as; \\
	$\pm$					& plus or minus; \\
	
	$y_{t}$, $y{k}$           	& output vector at continuous-time sampled instants;\\
	$z(t)$						& measurement output; \\
	$H(k)$, $L(k)$, $C$, $D$ 	& matrices of the observation model; \\
	$v(t)$, $v(k)$				& measurement noise vector; \\
	$t_i$ 					& continuous-time sampled index; \\
	$k$						& discrete-time index of measurements;  \\
	$T$						& expected time interval; \\
	$d(k)$					& observation sampling time delays; \\
	$l$						& amount of different known delays; \\
	$\xi(t)$				& random variable that models packet dropout; \\
	$\gamma(t)$				& random variable that models uncertain observation; \\
	$C_i(t)$				& real time approximation from the $i^{th}$ computer clock; \\
	$a_i$					& $i^{th}$ clock drift; \\
	$b_i$					& $i^{th}$ clock offset; \\
	$a_{ij}$				& relative drift between $i^{th}$ and $j^{th}$ node; \\
	$b_{ij}$				& relative offset between $i^{th}$ and $j^{th}$ node; \\

	$\Delta$				& measured value interval between observations; \\
	$\delta_k$				& random time interval; \\
	
	$p(k)$					& Bernoulli distribution parameter; \\
	$\epsilon_k$			& deviation from expected time instants; \\
	$\sigma$				& standard deviation; \\
	$\kappa$				& shape parameter; \\
	$\theta$				& scale parameter; \\
	
	$\mu s$					& micro seconds; \\
	$E.E$					& energy efficiency; \\
	$Comp.$					& complexity; \\

\end{longtable}


\subsection*{Chapter 4}

\begin{longtable}{ll}
%	$\mathbb{N}$			& natural numbers; \\
%	$\mathbb{R}$			& real numbers; \\
%	$\mathbb{R}^n$			& $\mathbb{R}^{nx1}$ $n$-dimensional Euclidean space; \\
	
%	$\mathcal{N}$			& Gaussian distribution; \\
	
	$\triangleq$			& equals by definition; \\	
%	$>$						& greater than; \\
%	$\geq$					& greater than or equal to; \\
%	$\sim$					& is distributed as; \\
%	$\pm$					& plus or minus; \\
	$\operatorname{arg}  
	\underset{x}
	{\operatorname{max}} 
					f(x)$	& argument $x$ that maximizes function $f(x)$; \\ 
					
	$\rho(x|y)$				& conditional probability density function of $X$ given $Y$; \\
	$x_k$						& state vector; \\
	$\hat{x_k}$				& estimated state vector; \\
	$y$						& observation vector; \\
	
%
%	$\sigma$				& standard deviation; \\
%	$\kappa$				& shape parameter; \\
%	$\theta$				& scale parameter; \\
	
	
\end{longtable}



\subsection*{Chapter 5}

\newpage
\section*{Acronyms}
\begin{longtable}{ll}
	CDF			& Cumulative Distribution Function; \\
	CU			& Covariance Union; \\
	DAI			& Data In; \\
	DAO			& Data Out; \\
	DEI			& Decision In; \\
	DEO 		& Decision Out; \\
	DMTS		& Delay Measurement Time Synchronization; \\
	DSET 		& Dempster-Shafer Evidence Theory; \\		
	EKF			& Extended Kalman Filter; \\
	FEI			& Feature In; \\
	FEO			& Feature Out; \\
	FIS			& Fuzzy Inference System; \\
	FISST		& Finite-Set Statistics; \\
	FTSP		& Flooding Time Synchronization Protocol; \\
	GPS			& Global Positioning System; \\
	HMM			& Hidden Markov Model; \\
	KF 			& Kalman Filter; \\
	UKF			& Unscented Kalman Filter; \\
	LEETS		& Lightweight and Energy Efficient Time Synchronization; \\
	LTS			& Lightweight Tree-based Synchronization; \\
	LS			& Least Squares; \\
	MAP			& Maximum A Posteriori; \\
	ML			& Maximum Likelihood; \\
	MMSE		& Minimum Mean Square Error; \\
	MOP 		& Measures Of Performance; \\
	NEES		& Normalized Estimation Error Squared; \\
	NIS			& Normalized Innovation Squared; \\
	NMR 		& Nuclear Magnetic Resonance; \\
	NTP			& Network Time Protocol; \\
	NUS 		& Non-Uniform Sampling; \\ 
	OOSM 		& Out-Of-Sequence Measurement; \\
	PDF         & Probability Density Function; \\
	PDAF		& Probabilistic Data Association Filter; \\
	PF 			& Particle Filter; \\
	RBS			& Reference Broadcast Synchronization; \\
	RMSE		& Root Mean Square Error; \\
	RTPF		& Real-Time Particle Filter; \\
	RV			& Random Variable; \\
	SOD			& Send-On-Delta; \\
	SNR			& Signal-to-Noise Ratio; \\
	TDS			& Time-Delay System; \\
	TDP			& Time-Diffusion Protocol; \\
	TPSN		& Timing-sync Protocol for Sensor Networks; \\
	TSST		& Time Synchronization based on Spanning Tree; \\
	TTF 		& Track-to-Track Fusion; \\
	UKF			& Unscented Kalman Filter; \\
	WSN			& Wireless Sensor Networks; \\
	ZOH			& Zero-Order Holder; \\
\end{longtable}


%===================================================================================================


%===================================================================================================
%                                         SUMÁRIO
%===================================================================================================
%!!!!!!!!!!!!!!!!!!!!!!!!!!!!!!!!!!!!!!!     ATENÇÃO     !!!!!!!!!!!!!!!!!!!!!!!!!!!!!!!!!!!!!!!!!!!
%	Para saber o número da primeira página do sumário, verifique o número da última página
% colocado na tabela de símbolos: arquivo simbolos.tex
%!!!!!!!!!!!!!!!!!!!!!!!!!!!!!!!!!!!!!!!!!!!!!!!!!!!!!!!!!!!!!!!!!!!!!!!!!!!!!!!!!!!!!!!!!!!!!!!!!!!
%%% Força a contagem da página.
\clearpage
\setcounter{page}{25}
%%%

\phantomsection
\addcontentsline{toc}{chapter}{SUMÁRIO}
\DeclareRobustCommand\contentsname
{\begin{center} \chapter*{SUMÁRIO} \end{center}}

\tableofcontents
\clearpage
%===================================================================================================


%===================================================================================================
%                                    DEMAIS FORMATAÇÕES
%===================================================================================================
%%% Setspace
\onehalfspacing         % Espaçamento de 1.5
%%%

%%% Iniciar capítulos apenas em páginas impares
\let\originalchapter=\chapter
\def\chapter{\cleardoublepage\originalchapter}
%%%

%%% Muda a numeração de páginas para formato arábico
\pagenumbering{arabic}
\setlength{\parskip}{1\baselineskip}
%%%
%==================================================================================================
